\section{Between Castle, Abbey, and New University}

\subsection{A birth without a birthday}

Thomas was born about 1225, probably at the Aquino family's castle of Roccasecca
in the Kingdom of Sicily.  Modern reconstructions ordinarily place the birth in
1224 or 1225, while a current alternative favors 1226.  No contemporary birth
act fixes the year, day, or room.  The range is reconstructed from later
statements about his age and from the sequence of his education.  ``Aquinas''
identifies the family and territorial association with Aquino; it is not a
surname in the modern sense.
His father Landulf belonged to the lesser nobility attached to the political
world of Emperor Frederick~II.  His mother Theodora came from the related
aristocratic networks of southern Italy.  Documentary traces of siblings and
property make this household more than hagiographic scenery, but they do not
yield the intimate childhood sometimes supplied by later retelling.

The region stood between papal and imperial power.  Monte Cassino was at once a
great Benedictine monastery, a landholder, a school, and a strategic fortress.
Family members held or sought ecclesiastical positions; placing a younger son
with the abbey could serve devotion, education, and lineage together.  Early
biographers say that Thomas was offered to Monte Cassino as a child.  The claim
is credible within the family's documented relations, but their descriptions
were written after his sanctity had become the subject.  They cannot establish
that anyone already foresaw a theologian or that an abbacy was a settled plan.

At the abbey Thomas encountered the Psalter, Latin grammar, liturgical time,
common discipline, and the material order of books.  A later question placed in
the child's mouth---what is God?---works perfectly as a hagiographic overture.
It may preserve remembered curiosity; it cannot be heard as a stenographic
record of vocation.  The historically stronger point is institutional: before
Thomas learned the rapid exchange of university disputation, he knew a house in
which study, prayer, property, patronage, and politics could not be separated.

\subsection{Naples and a changing curriculum}

Conflict around Frederick~II disrupted Monte Cassino, and Thomas continued his
education at the imperial studium in Naples, probably from about 1239.  Founded
to train servants without requiring travel to Bologna, the school joined civic
and royal purposes to a newly enlarged curriculum.  Latin translations made
more of Aristotle available; Arabic and Jewish commentators entered Christian
debate through translation; logic and natural philosophy pressed questions
that an inherited textbook order could not simply ignore.  Later sources name
teachers, especially Peter of Ireland and Martin of Dacia, but no transcript
reconstructs Thomas's exact course or first conclusions.

Naples also brought the Order of Preachers into view.  The Dominicans were a
young mendicant order whose common life, poverty, study, preaching, and mobile
obedience differed sharply from the stable and aristocratically entangled
future his family could imagine at Monte Cassino.  Thomas did not choose
``reason'' instead of monastic faith.  He chose one ecclesial form of religious
life over another, and the choice carried economic and dynastic consequences.
His later theology would insist that grace perfects rather than destroys
nature.  His own vocation first became public as a conflict over which natural
bonds and ecclesial goods should govern a son.

\evidenceaxis{Thomas studied at Monte Cassino and Naples before entering the
Order of Preachers.}{Family and ecclesiastical documents control the setting;
early Dominican lives and the canonization dossier supply the personal
sequence; modern chronology coordinates them.}{No school register gives a
complete curriculum, and later narratives interpret childhood from the saint
Thomas became.}
